\section{Parameters Description}

\subsection{Initial Bounds}
As mentioned in the introduction, our objective is to find parameters that make the signature ``better'' than SLH-DSA. Before we can do so, we have to fix the strict boundaries that we will not cross:
\begin{enumerate}
    \item We obviously do not want a signature larger than SLH-DSA - otherwise, why bother in the first place?
    \item For Bitcoin, we require at least a 128-bit security level. None of the parameters we select may degrade security below this floor.
    \item At least $q_{s} = 2^{40}$ supported signatures. SLH-DSA-128s defines the maximum number of signatures as $q_{s} = 2^{64}$, which is generous. For Bitcoin, we assume that $q_{s} = 2^{40}$ is a practical compromise, comfortably covering the foreseeable needs of on-chain operations and the Lightning Network.
    \item Sub-algorithms, hash, and addressing functions remain unchanged.
\end{enumerate}

\subsection{SLH-DSA Parameters}
SLH-DSA uses a hyper-tree of eXtended Merkle Signature Schemes (XMSS) to authenticate the underlying one-time and few-time signature keys. We formally parameterize each candidate configuration using the tuple $(h, d, k, a, w)$, where each variable controls a distinct geometric trade-off within the overall structure.

\paragraph{Hyper-Tree Geometry ($h$ and $d$).}
The hyper-tree configuration determines the capacity and layer-traversal costs of the signature scheme.
\begin{itemize}
    \item $h$ is the total height of the multi-layered hyper-tree. This dictates the maximum number of few-time signature instances supported at the base layer, yielding a total capacity of $2^h$ unique leaf nodes.
    \item $d$ is the number of independent tree layers comprising the hyper-tree.
\end{itemize}
To ensure a symmetric construction across all internal boundaries, our search pipeline requires that the total height $h$ be perfectly divisible by the layer count $d$. Consequently, each individual XMSS tree layer has a uniform height $h'$:
\begin{equation*}
    h' = \frac{h}{d}.
\end{equation*}
Adjusting the ratio between $h$ and $d$ governs the balance between signature size and key-generation cost. Reducing the number of layers $d$ shrinks the total number of intermediate public keys and authentication paths in the signature payload, saving valuable bytes. However, this exponentially increases the height $h'$ of each constituent tree, requiring significantly more hash operations to compute the root node during key generation.

\paragraph{Few-Time Signature Layer ($k$ and $a$).}
The lowest layer of the hyper-tree authenticates the Forest of Random Subsets (FORS) few-time signature scheme, which directly signs the message digest.
\begin{itemize}
    \item $k$ is the number of independent, parallel Merkle trees in the FORS layer.
    \item $a$ is the height of each individual FORS tree, so that each constituent tree contains exactly $t = 2^a$ leaves.
\end{itemize}

Increasing $k$ and $a$ strengthens the scheme against multi-target subset-forgery attacks and allows the structure to support larger capacities safely. However, expanding these dimensions directly increases the byte footprint of the appended authentication paths.

\paragraph{One-Time Signature Base ($w$).}
Internal tree roots are signed using the Winternitz One-Time Signature Plus (WOTS+) scheme. The base parameter $w$ defines the digit representation used to split the message digest during signing. Our framework evaluates candidate configurations across the standard baseline settings $w \in \{16, 32, 256\}$.

Selecting a larger base reduces the number of parallel hash chains $\ell$, which in turn reduces the overall WOTS+ payload size. This spatial efficiency comes at the cost of local compute: larger base parameters require longer hash chains, resulting in higher execution costs during key and signature generation.

We sweep the ranges as follows:
\begin{equation*}
    (h, d, k, a, w) \in \{40, \dots, 50\} \times \{2, \dots, 24\} \times \{2, \dots, 25\} \times \{8, \dots, 20\} \times \{16, 32, 256\}.
\end{equation*}

\textit{The SLH-DSA-128s baseline tuple (63, 7, 14, 12, 16) sits outside this range and is included for comparison}